\documentclass{cumcmthesis}

\begin{document}

\section{实验一代码}
\begin{lstlisting}[language=C++]
#include <iostream>
#include <iomanip>
#include <string>
using namespace std;
int RULENUM = 15;//动物识别系统的规则库有条规则
int CAUSENUM = 5;//各条规则中的前提条件最多有个
int count1 = 0;//记录所选择的动物特征的个数
/*事实*/
string fact[] = { "", "有毛发", "有奶", "有羽毛", "会飞", "会下蛋", "吃肉", "有犬齿", "有爪", "眼盯前方", "有蹄", "嚼反刍", "黄褐色", "身上有暗斑点", "身上有黑色条纹", "有长脖子", "有长腿", "不会飞", "会游泳", "有黑白二色", "善飞", "哺乳动物", "鸟", "食肉动物", "蹄类动物", "金钱豹", "虎", "长颈鹿", "斑马", "鸵鸟", "企鹅", "信天翁" };
/*规则*/
int rule[][6] = { { 0, 0, 0, 0, 0, 0 }, { 21, 1, 0, 0, 0, 0 }, { 21, 2, 0, 0, 0, 0 }, { 22, 3, 0, 0, 0, 0 }, { 22, 4, 5, 0, 0, 0 }, { 23, 6, 0, 0, 0, 0 }, { 23, 7, 8, 9, 0, 0 }, { 24, 21, 10, 0, 0, 0 }, { 24, 21, 11, 0, 0, 0 }, { 25, 21, 23, 12, 13, 0 }, { 26, 21, 23, 12, 14, 0 }, { 27, 24, 15, 16, 13, 0 }, { 28, 24, 14, 0, 0, 0 }, { 29, 22, 15, 16, 17, 19 }, { 30, 22, 18, 17, 19, 0 }, { 31, 22, 20, 0, 0, 0 } };
int conditionlist[24];//所选择的动物特征
bool match(int a);           //函数声明
int inference();
int main()
{
	cout << "以下是一些动物的特征:" << endl;
	for (int i = 1; i <= 24; i++)
	{
		cout << i << ".";
		cout << setiosflags(ios::left) << setw(14) << fact[i] << "  ";
		if (i % 4 == 0)
		{
			cout << endl;
		}
	}
	int a;
	int k = 0;
	cout << "请选择动物的特征:" << endl;
	while ((cin >> a))                      //a为整形的，按任意字母可结束循环
	{
		conditionlist[k] = a;
		k++;
		count1++;
	}
	cout << "该动物名称为：" << endl;
	inference();
	system("pause");
	return 0;
}
/*知识匹配*/
bool match(int a)
{
	int i = 1;
	int j = 0;
	int flag = 0;
	int flag1 = 0;
	int num = 0;
	while ((flag == i - 1) && (flag1 != count1))
	{
		flag = 0;
		for (i = 1; i <= CAUSENUM; i++)
		{
			if (rule[a + num][i] == 0)
				break;
			for (j = 0; j < count1; j++)                     //所选择的动物特征的条数
			{
				if (rule[a + num][i] == conditionlist[j])//匹配上一个条件flag计数加一
				{
					flag++;
					flag1++;
					break;
				}
			}
		}
		num++;
	}
	if ((flag != i - 1) || (count1 == 0))
	{
		return false;
	}
	else
	{
		for (i = 0; i < num; i++)
		{
			for (j = 0; j < num; j++)
			{
				if ((i != j) && (rule[a + i][0] != rule[a + j][0]))
				{
					return false;
				}
			}
		}
		return true;
	}
}
/*推理*/
int inference()
{
	int i;
	int j = 0;
	for (i = 1; i <= RULENUM; i++)
	{
		if (match(i) == true)
		{
			cout << fact[rule[i][0]] << endl;
			return 0;
		}
	}
	cout << "未知动物" << endl;
	getchar();
	return 0;
}
\end{lstlisting}

\newpage

\section{实验二代码}
\begin{lstlisting}[language=C++]
#include<iostream>
#include<algorithm>
#include<string>
#include<vector>
#include<cmath>
#include<queue>
#define N 6 // 棋盘/迷宫 的阶数 
using namespace std;

class Node
{
public:
	int x, y; // 节点所在位置	
	int F, G, H; // G:从起点开始，沿着产的路径，移动到网格上指定方格的移动耗费。
	// H:从网格上那个方格移动到终点B的预估移动耗费，使用曼哈顿距离。 
	// F = G + H 
	Node(int a, int b) :x(a), y(b) {}

	// 重载操作符，使优先队列以F值大小为标准维持堆 
	bool operator < (const Node& a) const
	{
		return F > a.F;
	}
};

// 定义八个方向 
int dir[8][2] = { { -1, -1 }, { -1, 0 }, { -1, 1 }, { 0, -1 },
{ 0, 1 }, { 1, -1 }, { 1, 0 }, { 1, 1 } };
// 优先队列，就相当于open表 
priority_queue<Node>que;
// 棋盘
int qp[N][N] = { { 0, 0, 0, 0, 0, 0 },
{ 0, 1, 1, 0, 1, 1 },
{ 0, 0, 1, 0, 0, 0 },
{ 0, 0, 1, 1, 1, 0 },
{ 0, 1, 1, 0, 0, 0 },
{ 1, 1, 0, 0, 0, 0 } };
bool visit[N][N]; // 访问情况记录，close表 
int valF[N][N];   // 记录每个节点对应的F值
int path[N][N][2]; // 存储每个节点的父节点

int Manhuattan(int x, int y, int x1, int y1); // 计算曼哈顿距离 
bool NodeIsLegal(int x, int y, int xx, int yy); // 判断位置合法性
void A_start(int x0, int y0, int x1, int y1); // A*算法 
void PrintPath(int x1, int y1); // 打印路径

/* ----------------主函数------------------- */
int main()
{
	fill(visit[0], visit[0] + N * N, false); // 将visit数组赋初值false
	fill(valF[0], valF[0] + N * N, 0); // 初始化F全为0 
	fill(path[0][0], path[0][0] + N * N * 2, -1); // 路径同样赋初值-1 

	//  // 起点 // 终点
	int x0, y0, x1, y1;
	cout << "输入起点：";
	cin >> x0 >> y0;
	cout << "输入终点：";
	cin >> x1 >> y1;
	x0--; y0--; x1--; y1--;

	if (!NodeIsLegal(x0, y0, x0, y0))
	{
		cout << "非法起点！" << endl;
		return 0;
	}

	A_start(x0, y0, x1, y1);  // A*算法 
	PrintPath(x1, y1);        // 打印路径 
	system("pause");
	return 0;
}

/* ----------------自定义函数------------------ */
void A_start(int x0, int y0, int x1, int y1)
{
	// 初始化起点 
	Node node(x0, y0);
	node.G = 0;
	node.H = Manhuattan(x0, y0, x1, y1);
	node.F = node.G + node.H;
	valF[x0][y0] = node.F;
	// 起点加入open表 
	que.push(node);

	while (!que.empty())
	{
		Node node_top = que.top(); que.pop();
		visit[node_top.x][node_top.y] = true; // 访问该点，加入closed表 
		if (node_top.x == x1 && node_top.y == y1) // 到达终点 
			break;

		// 遍历node_top周围的8个位置 
		for (int i = 0; i < 8; i++)
		{
			Node node_next(node_top.x + dir[i][0], node_top.y + dir[i][1]); // 创建一个node_top周围的节点 
			// 该节点坐标合法 且 未加入close表 
			if (NodeIsLegal(node_next.x, node_next.y, node_top.x, node_top.y) && !visit[node_next.x][node_next.y])
			{
				// 计算从起点并经过node_top节点到达该节点所花费的代价 
				node_next.G = node_top.G + int(sqrt(pow(dir[i][0], 2) + pow(dir[i][1], 2)) * 10);
				// 计算该节点到终点的曼哈顿距离
				node_next.H = Manhuattan(node_next.x, node_next.y, x1, y1);
				// 从起点经过node_top和该节点到达终点的估计代价
				node_next.F = node_next.G + node_next.H;

				// node_next.F < valF[node_next.x][node_next.y] 说明找到了更优的路径，则进行更新
				// valF[node_next.x][node_next.y] == 0 说明该节点还未加入open表中，则加入 
				if (node_next.F < valF[node_next.x][node_next.y] || valF[node_next.x][node_next.y] == 0)
				{
					// 保存该节点的父节点 
					path[node_next.x][node_next.y][0] = node_top.x;
					path[node_next.x][node_next.y][1] = node_top.y;
					valF[node_next.x][node_next.y] = node_next.F; // 修改该节点对应的valF值 
					que.push(node_next); // 加入open表
				}
			}
		}
	}
}

void PrintPath(int x1, int y1)
{
	if (path[x1][y1][0] == -1 || path[x1][y1][1] == -1)
	{
		cout << "没有可行路径！" << endl;
		return;
	}
	int x = x1, y = y1;
	int a, b;
	while (x != -1 || y != -1)
	{
		qp[x][y] = 2; // 将可行路径上的节点赋值为2 
		a = path[x][y][0];
		b = path[x][y][1];
		x = a;
		y = b;
	}
	// □表示未经过的节点， █表示障碍物， ☆表示可行节点 
	string s[3] = { "□", "█", "☆" };
	for (int i = 0; i < N; i++)
	{
		for (int j = 0; j < N; j++)
			cout << s[qp[i][j]];
		cout << endl;
	}
}

int Manhuattan(int x, int y, int x1, int y1)
{
	return (abs(x - x1) + abs(y - y1)) * 10;
}

bool NodeIsLegal(int x, int y, int xx, int yy)
{
	if (x < 0 || x >= N || y < 0 || y >= N) return false; // 判断边界 
	if (qp[x][y] == 1) return false; // 判断障碍物 
	// 两节点成对角型且它们的公共相邻节点存在障碍物 
	if (x != xx && y != yy && (qp[x][yy] == 1 || qp[xx][y] == 1)) return false;
	return true;
}
\end{lstlisting}

\newpage

\section{实验三代码}
\begin{lstlisting}[language=Matlab]
%% 清空环境变量
clear all
clc

%% 训练集/测试集产生
load spectra_data.mat
% 随机产生训练集和测试集
temp = randperm(size(NIR,1));
% 训练集——50个样本
P_train = NIR(temp(1:50),:)';
T_train = octane(temp(1:50),:)';
% 测试集——10个样本
P_test = NIR(temp(51:end),:)';
T_test = octane(temp(51:end),:)';
N = size(P_test,2);

%% BP神经网络创建、训练及仿真测试(R2014a)

% 创建网络
net = feedforwardnet(9);
% 设置训练参数
net.trainParam.epochs = 1000;
net.trainParam.goal = 1e-3;
net.trainParam.lr = 0.01;
% 训练网络
net = train(net,P_train,T_train);
% 仿真测试
T_sim_bp = net(P_test);

%% RBF神经网络创建及仿真测试

% 创建网络
net = newrbe(P_train,T_train,0.3);
% 仿真测试
T_sim_rbf = sim(net,P_test);

%% 性能评价

% 相对误差error
error_bp = abs(T_sim_bp - T_test)./T_test;
error_rbf = abs(T_sim_rbf - T_test)./T_test;
% 决定系数R^2
R2_bp = (N * sum(T_sim_bp .* T_test) - sum(T_sim_bp) * sum(T_test))^2 / ((N * sum((T_sim_bp).^2) - (sum(T_sim_bp))^2) * (N * sum((T_test).^2) - (sum(T_test))^2)); 
R2_rbf = (N * sum(T_sim_rbf .* T_test) - sum(T_sim_rbf) * sum(T_test))^2 / ((N * sum((T_sim_rbf).^2) - (sum(T_sim_rbf))^2) * (N * sum((T_test).^2) - (sum(T_test))^2));
% 结果对比
result_bp = [T_test' T_sim_bp' T_sim_rbf' error_bp' error_rbf']

%% 绘图
figure
plot(1:N,T_test,'b:*',1:N,T_sim_bp,'r-o',1:N,T_sim_rbf,'k-.^')
legend('真实值','BP预测值','RBF预测值')
xlabel('预测样本')
ylabel('辛烷值')
string = {'测试集辛烷值含量预测结果对比(BP vs RBF)';['R^2=' num2str(R2_bp) '(BP)' '  R^2=' num2str(R2_rbf) '(RBF)']};
title(string)
\end{lstlisting}
\newpage

\section{实验四代码}
\begin{lstlisting}[language=Matlab]
%% 清空环境变量
clear all
clc

%% 导入数据
load citys_data.mat

%% 计算城市间相互距离
n = size(citys,1);
D = zeros(n,n);
for i = 1:n
    for j = 1:n
        if i ~= j
            D(i,j) = sqrt(sum((citys(i,:) - citys(j,:)).^2));
        else
            D(i,j) = 1e-4;      
        end
    end    
end

%% 初始化参数
m = 50;                              % 蚂蚁数量
alpha = 1;                           % 信息素重要程度因子
beta = 5;                            % 启发函数重要程度因子
rho = 0.1;                           % 信息素挥发因子
Q = 1;                               % 常系数
Eta = 1./D;                          % 启发函数
Tau = ones(n,n);                     % 信息素矩阵
Table = zeros(m,n);                  % 路径记录表
iter = 1;                            % 迭代次数初值
iter_max = 200;                      % 最大迭代次数 
Route_best = zeros(iter_max,n);      % 各代最佳路径       
Length_best = zeros(iter_max,1);     % 各代最佳路径的长度  
Length_ave = zeros(iter_max,1);      % 各代路径的平均长度  

%% 迭代寻找最佳路径
while iter <= iter_max
    % 随机产生各个蚂蚁的起点城市
      start = zeros(m,1);
      for i = 1:m
          temp = randperm(n);
          start(i) = temp(1);
      end
      Table(:,1) = start; 
      % 构建解空间
      citys_index = 1:n;
      % 逐个蚂蚁路径选择
      for i = 1:m
          % 逐个城市路径选择
         for j = 2:n
             tabu = Table(i,1:(j - 1));           % 已访问的城市集合(禁忌表)
             allow_index = ~ismember(citys_index,tabu);
             allow = citys_index(allow_index);  % 待访问的城市集合
             P = allow;
             % 计算城市间转移概率
             for k = 1:length(allow)
                 P(k) = Tau(tabu(end),allow(k))^alpha * Eta(tabu(end),allow(k))^beta;
             end
             P = P/sum(P);
             % 轮盘赌法选择下一个访问城市
             Pc = cumsum(P);     
            target_index = find(Pc >= rand); 
            target = allow(target_index(1));
            Table(i,j) = target;
         end
      end
      % 计算各个蚂蚁的路径距离
      Length = zeros(m,1);
      for i = 1:m
          Route = Table(i,:);
          for j = 1:(n - 1)
              Length(i) = Length(i) + D(Route(j),Route(j + 1));
          end
          Length(i) = Length(i) + D(Route(n),Route(1));
      end
      % 计算最短路径距离及平均距离
      if iter == 1
          [min_Length,min_index] = min(Length);
          Length_best(iter) = min_Length;  
          Length_ave(iter) = mean(Length);
          Route_best(iter,:) = Table(min_index,:);
      else
          [min_Length,min_index] = min(Length);
          Length_best(iter) = min(Length_best(iter - 1),min_Length);
          Length_ave(iter) = mean(Length);
          if Length_best(iter) == min_Length
              Route_best(iter,:) = Table(min_index,:);
          else
              Route_best(iter,:) = Route_best((iter-1),:);
          end
      end
      % 更新信息素
      Delta_Tau = zeros(n,n);
      % 逐个蚂蚁计算
      for i = 1:m
          % 逐个城市计算
          for j = 1:(n - 1)
              Delta_Tau(Table(i,j),Table(i,j+1)) = Delta_Tau(Table(i,j),Table(i,j+1)) + Q/Length(i);
          end
          Delta_Tau(Table(i,n),Table(i,1)) = Delta_Tau(Table(i,n),Table(i,1)) + Q/Length(i);
      end
      Tau = (1-rho) * Tau + Delta_Tau;
    % 迭代次数加1，清空路径记录表
    iter = iter + 1;
    Table = zeros(m,n);
end

%% 结果显示
[Shortest_Length,index] = min(Length_best);
Shortest_Route = Route_best(index,:);
disp(['最短距离:' num2str(Shortest_Length)]);
disp(['最短路径:' num2str([Shortest_Route Shortest_Route(1)])]);

%% 绘图
figure(1)
plot([citys(Shortest_Route,1);citys(Shortest_Route(1),1)],...
     [citys(Shortest_Route,2);citys(Shortest_Route(1),2)],'o-');
grid on
for i = 1:size(citys,1)
    text(citys(i,1),citys(i,2),['   ' num2str(i)]);
end
text(citys(Shortest_Route(1),1),citys(Shortest_Route(1),2),'       起点');
text(citys(Shortest_Route(end),1),citys(Shortest_Route(end),2),'       终点');
xlabel('城市位置横坐标')
ylabel('城市位置纵坐标')
title(['蚁群算法优化路径(最短距离:' num2str(Shortest_Length) ')'])
figure(2)
plot(1:iter_max,Length_best,'b',1:iter_max,Length_ave,'r:')
legend('最短距离','平均距离')
xlabel('迭代次数')
ylabel('距离')
title('各代最短距离与平均距离对比')
\end{lstlisting}

\end{document}

